\section{Scenarios on the frozen board}
\label{sec:ssrboard}

The transport and its dial are complete; point them at the real data.  This section
runs one scenario --- the forward down $\MacSsrScenMovePct\%$ ---
across the entire frozen SPY board, all eight expiries, under the
three named regimes, and reads the results with everything the
chapter has taught.  Each expiry is transported from its \emph{own}
stored fit, and the ATM value after the move is the exact transport:
the node's fitted curve read at $k=\mathcal{R}H$.

\begin{figure}[htbp]
\centering
\paperfig{\textwidth}{fig_ssr_scenario.pdf}
\caption{The frozen SPY board under a $-\MacSsrScenMovePct\%$
scenario.  (a)~ATM implied volatility by expiry (log time axis):
today (orange, dashed) and the three transported boards.  The two-day
node moves from \MacSsrScenShortTodayPct\% to
\MacSsrScenShortRzeroPct/\MacSsrScenShortReindexPct/\MacSsrScenShortRtwoPct\%
under $\mathcal{R}=0/1/2$ --- its skew is
$s_0=\MacSsrScenShortSkew$.  (b)~The $\mathcal{R}=0$ vs
$\mathcal{R}=2$ delta gap \eqref{eq:ssrgap} at each expiry's 25-delta
put: from \MacSsrScenGapShortPts\ delta points at two days to
\MacSsrScenGapHeroPts\ at the December 2026 node and
\MacSsrScenGapLongPts\ at the far end.}
\label{fig:ssrscenario}
\end{figure}

Walk \cref{fig:ssrscenario}(a) from right to left.  At the long end
the board barely notices the regime: the December 2027 skew is gentle
($s_0(T)=\MacSsrScenLongSkew$), and the whole fan is
\MacSsrScenLongAtmMovePts\ volatility points wide.  Moving left, the three
transported curves fan apart, slowly and then violently.  At the
two-day node, today's ATM reading of \MacSsrScenShortTodayPct\% is
unchanged under sticky-moneyness, becomes
\MacSsrScenShortReindexPct\% under sticky-strike, and
\MacSsrScenShortRtwoPct\% under sticky-local-vol --- a spread of
\MacSsrScenShortAtmMovePts\ volatility points on one node, from one
tick.  The whole
spread is one number's doing: the two-day skew is
$s_0=\MacSsrScenShortSkew$ --- a two-day smile prices a
$\MacSsrScenMovePct\%$ move as a catastrophe, and its put wing is
priced accordingly --- and every formula in this chapter reads the
smile a multiple of the move down that wing.  The linear law would
say \MacSsrScenShortLinearPts\ points; the read is wider because the
two-day smile is sharply curved, the bend of \cref{fig:ssrdial}(b)
grown large.

One observation about the same panel.  On the two shortest expiries
the reads at $k=H$ and $k=2H$ lie \emph{beyond the last quoted
strike} --- the two-day node's quotes stop less than three percent
from the money.  What the transport consumes there is the fitted
model's published wing: Chapter~5's stated completion, not a quoted
value.  A scenario engine cannot avoid this --- a five-percent
question to a two-day smile is a question about its wing --- but a
desk should know which of its scenario numbers stand on quotes and
which on convention.

Panel~(b) prices the same fan in the units of \cref{sec:ssrdelta}:
the $\mathcal{R}=0$ versus $\mathcal{R}=2$ delta gap
\eqref{eq:ssrgap}, evaluated at each expiry's 25-delta put --- the
standard liquid put benchmark; its strike solves
$\PhiN(d_+)=0.75$ on the node's own fitted smile --- with the smile's
local slope at that strike.  The gap runs from
\MacSsrScenGapShortPts\ delta points at two days, through
\MacSsrScenGapHeroPts\ at the running December 2026 node, to
\MacSsrScenGapLongPts\ at the end of the board.  The gap is
$2\phin(d_+)\sqrt{\tau}$ times the slope at the benchmark put, and the
two factors pull against each other: the slope falls from
\MacSsrScenSlopeShort\ at two days to \MacSsrScenSlopeHero\ at
December 2026 and \MacSsrScenSlopeLong\ at the far end, while
$\sqrt\tau$ grows.  The slope falls faster than $1/\sqrt\tau$ over
the first weeks and more slowly after, so the product rises through
the December node and barely eases beyond it.  There is no safe
corner: the choice of
dial is worth sixteen to twenty-two delta points on the benchmark put
at \emph{every} expiry of an ordinary index board.

One paragraph on validity, where it bites.  A transported smile is a
new smile, and nothing in \eqref{eq:ssrshift} re-checks the butterfly
or calendar conditions of Chapter~2: for small moves the transported
surface inherits validity to first order, but at scenario sizes the
outputs should be re-checked before anything downstream consumes them
--- and there is a matching theorem: an arbitrage-free dynamics cannot
move the smile by level shifts alone indefinitely
\citep{RogersTehranchi2010}, so every rule in this chapter is a local
scenario device, not a global model.  The one output that is
arbitrage-free by construction is the frozen-field \emph{reprice} of
\cref{sec:ssrfield}, because it prices an actual model rather than
transporting a curve; the reference implementation offers it on
request, on the field extracted from the published surface, and reports
the realized ratio $\widehat{\mathcal{R}}$ beside it.

\subsection{What a spot move moves}
\label{sec:ssrmoves}

A transport rule is only half a policy; the other half says what it
is applied to, and the reference implementation's answers are short.
A spot move never refits.  The move is read from the forward,
$H_T=\log(F_T^{\rm new}/F_T^{\rm old})$ per expiry, with the forward
following the spot multiplicatively under a yield or proportional
dividend convention (one $H$ for every expiry) and additively,
$\Delta F=\Delta S\,e^{rt}$, under a cash schedule (Chapter~6's
elasticity split, so $H_T$ then differs by expiry).  The fitted
parameters of every node are left untouched; what is drawn and
priced is the transport \eqref{eq:ssrshift} of the fitted curve,
under one dial for all expiries and all names --- sticky-strike by
default.  The node's quotes keep their strikes: their labels slide by
the quote rule and their volatilities do not move, whatever the dial.
The node's prior (Chapter~10) is transported from its own forward by
the same rule, and a local-volatility sheet (Chapter~4) is not
re-marched: its nodal variances stay and its strike axis is relabeled
by the grid rule \eqref{eq:ssrgrid}, so that the sheet and the smiles
move together.  A recalibration resets the move to zero at the
chain's spot.

The transport also runs backward, on the quotes themselves.  A quote
observed at spot $S_i$ a few minutes ago belongs, at the current spot
$S$, to a slightly different smile; rather than trim it
(Chapter~6, \cref{rem:fwdmask}), the reference implementation moves
it.  With $\delta_i=\log(F/F_i)$ the forward move since the quote and
$w_0$ the node's last fit, the quote's total variance is shifted by
the amount the fit itself would have moved at that strike,
\[
\Delta w_i=w_0\big(x+(\mathcal{R}-1)\delta_i\big)-w_0(x),
\qquad x=k_i+\mathcal{R}\log(F_i/F_{\rm fit}),
\]
and relabeled at $k_i-\delta_i$: under the default $\mathcal{R}=1$
a stale quote keeps its volatility and only its label moves, exactly
the quote rule.  Its age enters as uncertainty: a band widened, or a
weight reduced, by $a\sqrt{\text{minutes}}$ in volatility, with $a$
a stated constant ($3.6$ vol bp per root-minute at a $15\%$
volatility level).  A chain whose quotes share one spot and one time
stamp is untouched.

\subsection{The contract}
\label{sec:ssrcontract}

The chapter closes as the others do: three columns of decreasing
strength --- proved, measured, chosen.

\begin{table}[htbp]
\centering
\small
\caption{The chapter's contract.}
\label{tab:ssrcontract}
\begin{tabular}{p{0.31\textwidth}p{0.31\textwidth}p{0.30\textwidth}}
\toprule
\textbf{Proved} & \textbf{Measured} & \textbf{Convention}\\
\midrule
A snapshot does not identify dynamics: every dial value prices
today's quotes identically (\cref{rem:ssrident}). &
The fan of tomorrows on the frozen node: \MacSsrMarkSpreadBp\ vol bp
at one strike for a $\MacSsrFanMovePct\%$ move
(\cref{fig:ssrquestion}). &
The dial $\mathcal{R}$ itself: assumed, not estimated --- nothing in
this chapter regresses realized smile moves on skew.\\
\addlinespace
Shape preservation leaves one free slide; ATM response
$\mathcal{R}s_0H$ (\cref{prop:ssratm}); fixed-strike read
$\sigma_{\mathrm{old}}(k_0+(\mathcal{R}-1)H)$, exactly
\eqref{eq:ssrstrike}. &
The second-order bend: \MacSsrDialBendBp\ vol bp at $H=-6\%$ under
sticky-strike, \MacSsrDialBendTwoBp\ under $\mathcal{R}=2$
(\cref{fig:ssrdial}). &
The transport is shape-frozen: a slide in total variance, chosen
over a vertical level (equivalent at first order); one dial for
every expiry and name, sticky-strike by default.\\
\addlinespace
Total delta
$\PhiN(d_+)+\phin(d_+)\sqrt{\tau}(\mathcal{R}-1)\sigma_{\mathrm{old}}'(k_0)$;
regime gap $2\phin(d_+)\sqrt{\tau}\lvert\sigma_{\mathrm{old}}'(k_0)\rvert$
\eqref{eq:ssrdelta}--\eqref{eq:ssrgap}. &
\MacSsrDeltaGapPts\ delta points at the
\MacSsrDeltaMarkPutDelta-delta put, \MacSsrDeltaMaxPts\ at the peak
(\cref{fig:ssrdelta}); \MacSsrScenGapShortPts--\MacSsrScenGapHeroPts\
points across the board (\cref{fig:ssrscenario}). &
Deltas quoted in the forward (Black) convention, discount set aside;
$\tau=t$ on this board (Chapter~8's SPY finding).\\
\addlinespace
The half-rule $s_0=b/2$; frozen-field ATM response $2s_0H$; the
relabeling \eqref{eq:ssrell} and its expansion
\eqref{eq:ssrellexp}; the grid rule \eqref{eq:ssrgrid} generating the
dial. &
Half-rule ratio \MacSsrHalfRatio; ATM move \MacSsrHalfAtmMoveBp\ vs
predicted \MacSsrHalfAtmPredBp\ vol bp (\cref{fig:ssrhalf});
$\widehat{\mathcal{R}}$ flat at \MacSsrRepFlatShort\ for the straight
field, $2+2cH/b$ (\MacSsrRepBentMean, \MacSsrRepBentBigMean) for the
bent one (\cref{fig:ssrreprice}). &
The frozen field is an approximate premise: measured markets sit
below $\mathcal{R}=2$ \citep{Bergomi2016}; the wings of the
answer depend on field shape the quotes do not pin down.\\
\addlinespace
--- (this row is empirical by design) &
The board scenario: two-day ATM spread
\MacSsrScenShortRzeroPct--\MacSsrScenShortRtwoPct\%; short-expiry
reads beyond the last quote (\cref{fig:ssrscenario}). &
Scenario reads on the two shortest expiries consume the published
wing (Chapter~5's stated completion); transported surfaces are to
be re-checked before downstream use \citep{RogersTehranchi2010};
what a move transports, and the age constant of a synchronised quote
(\cref{sec:ssrmoves}).\\
\bottomrule
\end{tabular}
\end{table}

Step back to the book's thesis one more time.  \emph{Validity} in
this chapter is modest and structural: transports preserve
arbitrage-freeness only locally, and the one globally safe path is to
move the model, not the curve.  \emph{Information} is the chapter's
sharpest lesson, because here the information runs out: today's
surface determines every strike derivative and no spot derivative
(\cref{rem:ssrident}) --- and even the regime that derives its answer
inherits the field's under-determination in the wings.
\emph{Convention} therefore carries more weight here than anywhere
else in this
book: the dial is a stated choice, worth double-digit delta points,
that no snapshot can test.  What \emph{could} test it is time ---
watching how smiles actually move as spot moves, tick after tick, and
weighing each new observation against what the transported surface
expected.  Weighing new evidence against a moving prior is a filtering
problem, and it is the subject of the next chapter.
